% RESEARCH_DRAFT_SHA256: 2d51bee6bc986cccd3eabfac0215ff71b3f717b1649370d34a57b1670a11839c
% Layout and source-reference wrappers expanded for the web renderer.
\documentclass[11pt]{article}
\usepackage[margin=25mm]{geometry}
\usepackage{amsthm}
\usepackage{fontspec,unicode-math}
\setmainfont{Latin Modern Roman}
\setmathfont{Latin Modern Math}
\usepackage{amsmath,amssymb,mathtools,xurl,hyperref}
\hypersetup{colorlinks=true,urlcolor=blue,linkcolor=blue}
\setcounter{secnumdepth}{0}
\setlength{\parindent}{0pt}
\setlength{\parskip}{5pt}
\setlength{\emergencystretch}{3em}
\newcommand{\sref}[2]{\hyperlink{source:#1:#2}{[S#2]}}
\newcommand{\cataloguescope}[1]{\paragraph{Source record.}#1}
\newcommand{\floor}[1]{\left\lfloor#1\right\rfloor}
\newtheorem{proposition}{Proposition}
\newtheorem{lemma}[proposition]{Lemma}
\begin{document}
\paragraph{Research manuscript.} OpenAI \texttt{gpt-6-astra}, reasoning \texttt{ultra}; 11 October 2026.
\paragraph{Preserved source body.} The complete catalogue mathematical text and sources follow with layout and source links expanded, including source qualifications and its historical preparation notes. New research begins at the section headed Research attempt.

% UnsolvedMath ID 30002629; source catalogue code OWR-13105-008
\section{Depth-Independent Attribute-Based Encryption}\label{30002629}
{\small UnsolvedMath ID 30002629 · Computer Science}
\subsection{Definitions and mathematical statement}

\paragraph{Shared notation.}Standard number systems are $\mathbb N,\mathbb Z,\mathbb Q,\mathbb R,\mathbb C$. Specialized notation and index ranges follow the cited source and the definitions below.


Attribute-based encryption associates ciphertexts or keys with attributes and access policies, permitting decryption when the policy is satisfied. Circuit depth is the longest gate path in the represented policy; key-homomorphic encryption supports prescribed algebraic combinations of keys.
\paragraph{Record scope and qualifications.}
The catalogue title is retained as metadata. Named certificates, contracts, bounds or frameworks in that title are not asserted as results of the original source extract.
Which parameters must be depth-independent and the policy/security models are not specified in the supplied question.

Source-record limitation: Which encryption parameters must be independent of depth and the security model are absent; the additional applications question is an agenda.
\paragraph{Statement recorded in the supplied source.}
Can attribute-based encryption be constructed with parameters independent of circuit depth? What applications does fully key-homomorphic encryption have? \href{https://doi.org/10.4171/owr/2014/35}{[S2]}
\paragraph{Source status.}
The supplied catalogue labels this source record partially solved; that metadata was checked against the supplied snapshot on 2026-10-11. The label is preserved for provenance and is not an independent verification of present open status or mathematical correctness. The catalogue literature assessment is dated 2026-08-21.
\subsection{Short English statement}
Make attribute-based encryption independent of policy depth and identify applications of fully key-homomorphic encryption.
\subsection{Sources}
\begin{itemize}
\item[S1] Ulam AI and the UnsolvedMath Contributors, UnsolvedMath: A Curated Collection of Open Mathematics Problems, record 30002629 (OWR-13105-008). Catalogue curation and metadata: CC BY 4.0. The source extract is retained; definitions, English summaries and qualifications are editorial additions.
\url{https://huggingface.co/datasets/ulamai/UnsolvedMath}.
\item[S2] Cryptography (2015), Oberwolfach Reports, DOI 10.4171/owr/2014/35, pp. 1933--1990
\url{https://doi.org/10.4171/owr/2014/35}.
\item[S3] Eurocrypt 2014 fully key-homomorphic encryption
\url{https://www.iacr.org/archive/eurocrypt2014/84410298/84410298.pdf}.
\item[S4] Crypto 2013 identity and attribute-based FHE
\url{https://www.iacr.org/archive/crypto2013/80420297/80420297.pdf}.
\item[S5] Puncturable encryption from key-homomorphic encryption
\url{https://arxiv.org/abs/2007.06353}.
\item Source-specific Oberwolfach Reports licences, including CC BY-SA 4.0 where applicable, are retained; see THIRD\_PARTY\_NOTICES.md.
\end{itemize}
\paragraph{Preparation and review.}
The original source and UnsolvedMath contributors retain attribution for the source extract. Codex (GPT-6; exact serving revision not exposed) assisted with definitions, English summaries, scope notes and formatting. No new human mathematical authorship or independently verified solution is asserted. Permanent source ID 30002629 and code OWR-13105-008 checked against the complete supplied registry and current repository. Distinct upstream records, including repeated questions and broad or damaged extracts, are retained for complete coverage of the source CS category. A repeated record is not a new independent conjecture. Complete mathematical-equivalence and current-literature review has not been performed. Preparation and final review follow contribution prompts C and D, with source gaps disclosed. Structural and rendering checks do not establish a new theorem or open status.
\label{end:30002629}

\section{Research attempt}
\subsection{Literature, source context and scope}
The full original catalogue body and source qualifications are preserved above. The relevant primary workshop passages were read in the cached publisher report. Historical claims and subsequent attributed results are distinguished from the mathematical work below. A restricted example, standard lemma or known theorem is not counted as new progress on the entire original target. Missing assumptions and repeated or rhetorical extracts are not silently turned into different solved problems.

\subsection{Translation functionality alone does not establish an application}
The source's fully key-homomorphic encryption translates a ciphertext under attribute vector $x$ into one decryptable under the derived key label $(f(x),f)$, preserving its message. This is not merely additive homomorphism of ordinary encryption keys. We test the proposed route from that functionality to key-policy attribute-based encryption.
\begin{proposition}
Assume a translation scheme has correctness for every allowed function $f$. Giving a policy holder the secret key for label $(1,f)$ allows correct decryption of a ciphertext under $x$ whenever $f(x)=1$, by first translating it using $f$. However, these correctness statements alone do not imply secrecy of the resulting system.
\end{proposition}
\begin{proof}
On an encryption of $m$ under $x$, the translation correctness hypothesis says that applying the public operation for $f$ yields a ciphertext decryptable with the key labeled $(f(x),f)$ and returns $m$. If $f(x)=1$, this is exactly the supplied policy key, proving the conditional correctness statement.
For the second assertion, consider the deliberately insecure syntax $\operatorname{Enc}(x,m)=(x,m)$ and $T_f(x,m)=(f(x),f,m)$. Let keys be their public labels; decryption returns the embedded message when labels match and $\bot$ otherwise. This syntax satisfies the displayed functional translation and decryption requirements perfectly. Yet an observer reads $m$ directly from the initial or translated ciphertext, without any key. Thus functional correctness cannot by itself entail message secrecy or an attribute-based-encryption security game.
\end{proof}
The counterexample is not claimed to satisfy the source's security definition, compactness requirements, or computational assumptions; violating secrecy is its purpose. The source already derives applications using security as well as functionality. The failed attempted inference here is dropping that security premise when transferring an application from the translation identity alone.
\subsection{Verification and remaining gaps}
For two-bit attributes, three simple Boolean policies and both one-bit messages, every satisfied-policy translation was checked to decrypt correctly in the toy syntax, while direct reading recovers every message. This finite enumeration verifies the separation of correctness from confidentiality; it provides no secure primitive. No weaker assumption, greater supported circuit depth, shorter key, or new secure application is established. Those are the substantive source targets and remain unresolved in this attempt.

\section{References and source attribution}
Primary workshop report: \url{https://doi.org/10.4171/owr/2014/35}. Relevant passages read: 2014/35 pp. 1985--1986, fully key-homomorphic ciphertext translation, compactness, depth and applications.

All retrievals were checked on 11 October 2026. Original question attribution remains in the preserved Sources section. External theorems are attributed and not claimed as new discoveries or fully verified merely by citation.


\paragraph{Integration review.} Codex (GPT-6; exact serving revision not exposed) reviewed the written arguments and expanded layout and source-reference wrappers. Repository metadata and the production web renderer were checked using Python 3.13 and JavaScript with MathJax 3.2.2. This model-assisted review is not independent human proof verification; no PDF compilation or formal proof checking is asserted.

\section{Completion Estimate}
\noindent\textbf{COMPLETION: 0\%}
No substantiated new progress toward the entire original target is claimed. The worked results and failed route are delimited above; a known result, illustrative model, or auxiliary obstruction is not counted as original-target completion.

\section{Final status and attribution}
\textbf{Attempt status: unresolved.} This describes this research attempt, not a new ruling on the source catalogue's status or on external literature claims.
The mathematics and this self-audit were generated by OpenAI \texttt{gpt-6-astra}, reasoning \texttt{ultra}, on 11 October 2026. Codex (GPT-6; exact serving revision not exposed) orchestrated the batch, reviewed the written mathematical arguments, and contributed formatting/repository review. No human mathematical authorship, independent human proof checking, or formal verification is claimed. The historical preparation paragraphs in the reproduced catalogue body concern that body only.

Contribution prompt A was applied to the complete definition and supplied source record, followed by the adversarial checks of prompt B. Local reproduction files are \texttt{batch019/checks.py}, \texttt{computations.json}, and the per-ID audit JSON. They were not public links at preparation time. Numerical checks are finite and do not prove unrestricted results. No LaTeX engine was found locally; compilation was not performed by this research agent.

\end{document}
